% ============================================================
%  文档类选项（在 main.tex 的 \documentclass[...]{tongjithesis} 中设置）
%  field=science：理工科编号格式 1 / 1.1 / 1.1.1
% ============================================================

% ============================================================
%  封面信息
% ============================================================
\school{电子与信息工程学院}
\major{微电子科学与工程}
\student{2253644}{云发鹏}
\thesistitle{基于电子对抗网络的空地协同部署优化}{}
\thesistitleeng{Air-Ground Cooperative Deployment Optimization Based on Electronic Countermeasure Networks}{}
\thesisadvisor{徐凡}
\thesisdate{2026}{5}{15}

% ============================================================
%  信息说明页
% ============================================================
\infotype{design}

\infoabstract{%
本文研究复杂区域中的空地协同雷达/电子对抗节点部署问题。方法基于MOPSO-DT，引入Hertel-Mehlhorn凸分解、垂直交点坐标变换和A2G/G2G异构传播评估，以ECR和$J_{\min}$描述覆盖-压制取舍。实验在雷达方程、挑战性区域和调优场景中，与原始MOPSO-DT、NSGA-II、MOEA/D、SPEA2及随机搜索同预算比较。结果表明，本文流程能输出可解释的非支配部署方案，但经典MOEA方法在部分Pareto质量指标上仍有优势。
}

\infodrawings{19}
% 毕业设计：主要图纸张数与正文总字数。最终编译后按统计结果更新。
\infowordcount{16884}

\infomaterials{
  \item 项目程序源代码；
  \item 实验脚本、测试用例与辅助工具；
  \item 实验结果、生成图表、验证报告与复现说明文档。
}

